\section{Measure root event boundary package}
\label{sec:measure-root-event-boundary-package}

The root event boundary is the public face of the $\MeasureUp$ root threshold
seen by the downstream analysis chapters. It keeps the boundary over
$\mathsf{BHist}$ presentations, componentwise $\hsame$ event comparison,
$\NameCert_{\RealUp}$ endpoint comparison, $\StreamNameUp$ countable names,
and the displayed measure ledger rows.

\begin{theorem}[Measure root event-boundary carrier obligation]
\label{thm:measure-root-event-boundary-carrier}
The $\MeasureUp$ root event boundary exposes exactly the carried
$\mathsf{BHist}$ measure code, visible event rows, empty-event row, finite
union rows, countable union rows, relative-difference row, $\RealUp$ endpoint
row, and countable-sum witnesses supplied by
\autoref{def:measureup-carrier-surface},
\autoref{def:measure-root-threshold-obligation-package}, and
\autoref{thm:measure-root-threshold-carrier-classifier-obligation}. Every
exposed row is a displayed $\mathsf{BHist}$ row or a named endpoint row, and
the boundary adds no host set equality, quotient event identity, analytic
completion object, or standard bridge row.
\end{theorem}
\begin{proof}
Start from the carrier surface of \autoref{def:measureup-carrier-surface}.
It names the carried measure code, the measurable-event family, visible
event rows, empty-event row, finite and countable union rows,
relative-difference row, nonnegative endpoint row, and countable-sum
witnesses.

Next read the same rows through
\autoref{def:measure-root-threshold-obligation-package}. The threshold
package places them at the root over $\mathsf{BHist}$ presentations,
componentwise $\hsame$, $\NameCert_{\RealUp}$ endpoint comparison, and
$\StreamNameUp$ countable names.

Apply \autoref{thm:measure-root-threshold-carrier-classifier-obligation}.
It repacks the carrier and visible ledger rows as root boundary data and
keeps endpoint movement inside the $\RealUp$ classifier.

The exactness clauses of the carrier and threshold packages identify the
displayed ledger as the whole boundary surface. Therefore the event boundary
contains precisely the listed $\mathsf{BHist}$ rows and endpoint witnesses,
with no additional host equality or quotient event row.
\end{proof}

\begin{theorem}[Measure root event-boundary classifier obligation]
\label{thm:measure-root-event-boundary-classifier}
The public classifier on the $\MeasureUp$ root event boundary is
componentwise $\hsame$ on the visible event rows together with
$\NameCert_{\RealUp}$ classifier agreement on endpoint rows. Equivalently,
the boundary classifier is the root-facing instance of
\autoref{def:measure-classifier-specification} projected through
\autoref{thm:measure-root-threshold-carrier-classifier-obligation}: it
compares displayed event rows and $\RealUp$ endpoints only, and it never
asserts host set equality or quotient event identity.
\end{theorem}
\begin{proof}
Project the public classifier from
\autoref{def:measure-classifier-specification}. Its event component reads
visible event presentations through $\hsame$, while its endpoint component
reads measure values through $\NameCert_{\RealUp}$ and the $\RealUp$
classifier.

Now apply
\autoref{thm:measure-root-threshold-carrier-classifier-obligation}. The
root threshold projection restricts the classifier to the displayed
$\mathsf{BHist}$ event rows, ledger rows, and endpoint rows carried by the
boundary.

It remains to identify what the projection excludes. Since the carrier
surface and the threshold package list only $\mathsf{BHist}$ rows,
componentwise $\hsame$, $\NameCert_{\RealUp}$ comparison, and displayed
ledger witnesses, the classifier has no field in which host set equality or
quotient event identity could be read.
\end{proof}

\begin{theorem}[Measure root event-boundary downstream unblock]
\label{thm:measure-root-event-boundary-unblock}
The same $\MeasureUp$ root event boundary supplies the downstream rows
consumed by $\IntegralUp$, $\ProbSpaceUp$, and $\FunctionalAnalysisUp$. The
unblocking rows are exactly the countable-additivity and analysis-consumer
obligations of
\autoref{thm:measure-root-threshold-sigma-additive-obligation} and
\autoref{thm:measure-root-threshold-analysis-consumer-obligation}, read
through \autoref{thm:measure-integral-consumer-rows},
\autoref{thm:measure-probability-consumer-rows}, and
\autoref{thm:measure-functionalanalysis-consumer-rows}. They remain scoped to
$\mathsf{BHist}$, $\hsame$, $\NameCert_{\RealUp}$, $\StreamNameUp$, and the
displayed $\MeasureUp$ ledger rows.
\end{theorem}
\begin{proof}
First build the additive part of the boundary. Apply
\autoref{thm:measure-root-threshold-sigma-additive-obligation}. It supplies
the empty-zero row, finite disjoint-union readback, relative-difference
additivity, countable disjoint-union comparison, head-tail readback, and
countable-sum classifier transport over the same $\mathsf{BHist}$ rows.

Next build the consumer projection. Apply
\autoref{thm:measure-root-threshold-analysis-consumer-obligation}. It reads
the $\IntegralUp$, $\ProbSpaceUp$, and $\FunctionalAnalysisUp$ outgoing rows
from the root package without adding a host carrier, quotient event, analytic
completion object, or standard bridge row.

For the integral consumer, use
\autoref{thm:measure-integral-consumer-rows}. The consumed rows are the
covered measurable-event rows, $\RealUp$ endpoint classifier rows,
countable-sum rows, head-tail readback, and finite disjoint-union readback.

For the probability consumer, use
\autoref{thm:measure-probability-consumer-rows}. The consumed rows are the
total-event normalization row, empty-zero row, and finite disjoint-union
readback row.

For the functional-analysis consumer, use
\autoref{thm:measure-functionalanalysis-consumer-rows}. The consumed rows
are the nonnegative endpoint, monotone-inclusion, relative-difference, and
countable-sum classifier rows.

All three consumer projections point back to the same root event boundary.
Thus the boundary supplies the three downstream faces through countable
additivity and analysis-consumer rows while staying inside
$\mathsf{BHist}$, $\hsame$, $\NameCert_{\RealUp}$, $\StreamNameUp$, and the
displayed $\MeasureUp$ ledger.
\end{proof}
